\documentclass[11pt]{article}
\usepackage[margin=1in]{geometry}
\usepackage[T1]{fontenc}
\usepackage{amsmath,amssymb,array,longtable,xurl,hyperref}
\setlength{\emergencystretch}{2em}
\title{TVSC Workflow Strategy and Session Handoff}
\author{}
\date{September 15, 2026}
\begin{document}
\maketitle

\section{Read this first after resuming}
This is an internal strategy and decision record, not a third research
paper. It records the primary-documentation review conducted on September
15, 2026 and the next implementation steps. Statements marked
\emph{agreed} record user decisions; statements marked \emph{proposed}
are recommendations awaiting confirmation. Documentation findings are
cited separately from design choices. A first raw preparation implementation
has now been written, and the user supplied a passing R test transcript on
September 15, 2026. This is an externally reported run, not an independent
local reproduction. No completed package, runtime benchmarks or new
California estimates are reported here.

\paragraph{Current implementation checkpoint.}
The September 15 design discussion supersedes the earlier implementation
sequence. The user objected to advancing metrics and numerical construction
without first discussing their contracts. The \path{tvsc-problem.R.txt}
draft and its fixtures are set aside, not deleted or validated. It remains
a dense primary-problem builder, not a solver. Do not resume solver work
or interpret passing fixtures as approval of an interface or calibration
method. The approved first refactor is now written: preparation and its
raw blocks use version 2, and splitting accepts revised preparation while
retaining version-1 output. The user-supplied September 15, 2026 R console
transcript shows both revised fixture suites passing without unexpected
errors. It includes source-loading commands and a passing isolated-environment
preparation check without a public \texttt{tvsc.blocks()} dependency.
Session information and source revision identifiers were not supplied;
a clean-start session is not established. Independent local reproduction
remains pending because R is unavailable in this workspace. Metric
construction is now written as a separate source draft with new fixtures.
A subsequent September 15, 2026 user-supplied excerpt shows preparation and
metric construction sourced into an isolated environment, a passing
constant-mode comparison with the existing reference object, and the
suite success message. No local split or standalone scaling binding is
present. This supports a reported suite pass, but only the final check is
shown; earlier assertions and reference creation are not established by
the excerpt. The next validation checkpoint is a complete clean-session
record of all three revised suites, with session information and source
identifiers, not solver work. Independent reproduction remains pending.
No GitHub integration summary or independent test
record has been received for the preceding handoff.

\paragraph{Agreed revision: September 15, 2026.}
Section~\ref{sec:calibration-revision} records the current decisions.
The historical source contracts below remain descriptions of existing
drafts, not the final public interface. After design approval, the first
preparation/split refactor changed their source and test files. In
particular, \texttt{tvsc.vmetrics()} now implements donor-only scaling and
calibration in its separate source draft, with a reported user-run pass
supported by the final test excerpt and independent reproduction pending.
The old standalone block, scaling, metric and problem drafts remain
version-1 components and do not accept revised preparation/raw blocks.
Their earlier suites are not a runnable version-2 integration suite.

On September 15, 2026, the user approved one consistent \texttt{tvsc.*}
prefix for preparation, splitting, blocks, scaling and metrics. The five
source files and five fixture files now use \texttt{tvsc-} filenames;
definitions, internal calls and documentation follow the new names.
No compatibility aliases are retained. Arguments, defaults, formulas,
return classes and schema versions are unchanged. Current names are used
in the chronology below, but earlier passing transcripts used the former
\texttt{panel.*} names. A subsequent user-supplied transcript on September
15, 2026 includes all five renamed function definitions and fixture
creation. Preparation, splitting, blocks and scaling reach their success
messages. Metrics passes the displayed checks through dimensions, then
stops at a fixture's incorrect expected-error message. Assigning null
through the scaling field removes that list member, so recipe validation
fails before the expected scaling-metadata check. The fixture now tests
the absent field separately and retains a null-valued member for the
later check. No function-body change is required. The user subsequently
reports a metrics pass on September 15, 2026, supplying only the final
misaligned-divisor-names assertion and ``All panel metrics checks passed.''
Record this as a user-reported pass, not review of a complete rerun:
the excerpt alone does not establish execution of all earlier assertions.
The full rerun transcript, source-loading commands, revisions, session information and
evidence of a clean start are absent. Independent reproduction remains
pending for all five functions.

\path{R/tvsc.dataprep.R} contains the first raw-only
\texttt{tvsc.dataprep()} implementation. \path{tests/test-tvsc.dataprep.R}
contains executable fixtures whose checks passed in the user-supplied
console transcript. The expanding-window split constructor is now written
in \path{R/tvsc.split.R}, with fixtures in
\path{tests/test-tvsc.split.R}. A second user-supplied console transcript on
September 15, 2026 shows those split checks completing without errors and
ends with ``All chronological panel split checks passed.'' It includes the
eight-date example, stride/tail boundaries, period spacing, all-unit date
assignments, input preservation and invalid-input checks. Its source-loading
commands, exact function revisions and R session information were not
included. Neither test run was independently reproduced here.
The raw contemporaneous builder is now written in \path{R/tvsc.blocks.R},
with fixtures in \path{tests/test-tvsc.blocks.R}. A third user-supplied R
transcript on September 15, 2026 shows all displayed block assertions
completing without errors, ending with ``All raw predictor block checks
passed.'' This includes matrix values and orientation, naming/order,
cutoff isolation, single-predictor dimensions, outcome exclusion, pre-scaled
predictor preservation, transformed-outcome passthrough and invalid inputs.
Source-loading commands, initial panel creation, exact revisions and R
session information are absent. Record this as user-run evidence, not as
a locally reproduced or confirmed clean-session run. The scaling function
is now written in \path{R/tvsc.scale.R}, with fixtures in
\path{tests/test-tvsc.scale.R}. A fourth user-supplied R transcript on
September 15, 2026 includes fixture creation and shows all displayed
scaling assertions completing without errors, ending with ``All panel
scaling checks passed.'' It covers the four methods, zero-scale fallbacks,
outcome preservation, common-shift identities, cutoff isolation, metadata,
single-predictor dimensions, repeated-scaling rejection and numerical/input
failures. Source-loading commands, exact revisions and R session information
are absent. Independent reproduction and the standalone example remain
pending; do not infer cross-version or exhaustive validation. The next
source draft, \path{R/tvsc.metrics.R}, now constructs uniform metrics
and validates supplied common vectors or date-varying matrices. Its
fixtures are in \path{tests/test-tvsc.metrics.R}; the initial failure, fixture
correction and subsequent reported pass are recorded above. Neither the corrected suite nor
its standalone example has been executed in R here. There is no
empirical calibration, fitting, prediction or ATT implementation yet.
The files use \texttt{.R.txt} extensions because this Prism project
tracks text files but excludes \texttt{.R} files. R can source these text
files directly. The transcript shows assertions completing without errors
and the final message ``All raw panel preparation checks passed.'' The
source-loading command, exact source revision and R session information
were not included. Record those details and execute the separate documented
example before publication; do not infer cross-version compatibility or
exhaustive correctness from this test result.

\paragraph{Agreed no-additional-scaling option: September 15, 2026.}
The researcher may supply predictors that were transformed beforehand.
``Raw'' throughout the current preparation and assembly interface means
as supplied, not necessarily original economic units. Both existing
functions already leave those values unchanged. The future scaling
interface must support \texttt{scaling = "none"} as an explicit identity
operation, not as an implicit default. This option is agreed but its
interface is not implemented; do not add it to existing function calls
or replace their \texttt{recipe\$scaling = NULL} metadata manually.
That metadata describes package-applied processing only.

Record upstream transformations, affected columns, common target/donor
parameters, units and estimation samples separately until a provenance
interface is designed. Learned parameters must respect every validation
cutoff; slicing full-period-standardized data cannot undo prior leakage.
Fixed external unit conversions do not require full-panel estimation.
Retain the outcome in reporting units and use a separate scaled-outcome
predictor when desired. Current functions perform no inverse transformation
or upstream provenance certification. The block fixtures now include
pre-scaled predictor preservation and transformed-outcome passthrough;
these additions passed in the third user-supplied transcript, with a
clean-session reproduction and environment record still pending.

\paragraph{Function documentation convention.}
Each implemented public function should carry a self-contained
\texttt{\#'} header describing its purpose, arguments, returned object,
limitations and a small example, together with short \texttt{\#} comments
at substantive steps in the function body. This follows the user's requested
source-file documentation style. The convention is now applied to
\path{R/tvsc.dataprep.R}; future functions should follow it when implemented.
The repository overview remains in \path{README.md}. Documentation additions
do not establish runtime validation or package readiness.

\paragraph{Incremental GitHub publication.}
The user wants each function uploaded when it is finished, rather than
waiting for the entire software workflow. For each function, complete the
implementation, source-file documentation, executable example and R tests;
review actual test results and resolve failures before marking it ready.
Then publish to the agreed GitHub destination and record the published
revision. Bring any externally made fixes back into this workspace before
building dependent functions. Written, tested and uploaded are separate
statuses: preparation, splitting, block assembly and scaling are written and documented,
each with a user-supplied passing test transcript. Independent local execution
and GitHub publication remain unconfirmed.

\paragraph{Latest decisions: September 15, 2026.}
The user has approved separating preparation from chronological splitting.
\texttt{tvsc.dataprep()} requires no training/validation split;
\texttt{tvsc.split()} is a separate optional operation. A fixed-penalty
fit can use the full eligible pre-period without splits. Chronological
penalty selection still requires them and is not merely an RMSE reporting
exercise. The scaling scope is now limited to \texttt{sd} (default),
\texttt{zscore}, \texttt{none} and \texttt{custom}, with the panel-specific
formulas below. Pooled-level and robust alternatives are deferred.
This paragraph and the revised contracts below supersede the earlier
three-function-only workflow and pooling proposal.

\begin{itemize}
\item Paper one is \path{tvsc-paper.tex}, compiled through
\path{main.tex}. It contains the conditional theory and proofs and will
eventually report the empirical application. Statistical consistency is
explicitly out of the current scope.
\item Paper two is \path{tvsc-workflow-paper.tex}. It is a software-design
draft containing the chronological workflow and proposed interfaces.
\item Work order: use paper one's reviewed conditional theory as the
working baseline; specify, implement and validate the package for paper
two; return to paper one for the California analysis.
\item Preparation, splitting and block assembly now have user-supplied passing test
transcripts. The user has agreed to separate raw contemporaneous block
assembly from scaling. Complete their publication records, then discuss the
scaling interface with the agreed no-additional-scaling option. The default
and additional scaling rules remain open. Preserve all three interfaces.
\item Neither \texttt{R} nor \texttt{Rscript} was discoverable on the
current shell's PATH during this review. Recheck on resumption. Do not
claim R tests have run here. The user-supplied transcript is external
execution evidence; obtain session information and the tested source
revision for the publication record. No dependencies were installed.
\end{itemize}

\section{Agreed methodological and interface decisions}
The fitted objective is
\[
 (1-\lambda)\frac1T\sum_{t=1}^T
 \|X_{1,t}-X_{0,t}w_t\|_{V_t}^2
 +\frac{\lambda}{T-q}\sum_{t=q+1}^T\|\Delta^q w_t\|_2^2,
 \qquad w_t\in\Delta_J,
\]
with $q\in\{1,2\}$, $0\leq\lambda<1$, and predetermined diagonal
metrics whose entries are nonnegative and sum to one at every date.
Predictor weights are never estimated jointly or alternately with donor
weights. Nonunique primary optima use minimum-norm secondary selection,
not an added ridge penalty.

Agreed defaults are $q=1$, uniform predictor weights $v_{k,t}=1/K$,
chronological selection from 100 values
$\{0,0.01,\ldots,0.99\}$, and projected increment continuation.
The equivalent additive coefficient is $\alpha=\lambda/(1-\lambda)$.
The grid covers $\alpha$ from zero to 99, not all possible penalties.
The default order is not presumed to be the empirical winner.

The agreed separation of responsibilities gives the following intended
sequence for chronological selection; the split constructor requires
explicit initial and horizon sizes, with a default stride of one date:
\begin{verbatim}
prepared <- tvsc.dataprep(data, ...)
splits <- tvsc.split(prepared, initial = ..., horizon = ..., step = ...)
fit <- tvsc.fit(prepared, q = 1, predictor_weights = "uniform",
                lambda = "cv", splits = splits,
                lambda_grid = seq(0, 0.99, by = 0.01),
                cv_continuation = "projected_increment")
effects <- tvsc.att(fit, prepared,
                    continuation = "projected_increment")
\end{verbatim}
A fixed-penalty alternative needs no split object:
\begin{verbatim}
fit <- tvsc.fit(prepared, lambda = 0.4)
\end{verbatim}
Here 0.4 is an illustrative input, not a selected or recommended default.
The prepared panel is the same in both workflows. Direct fitting provides
in-sample diagnostics, not held-out predictive accuracy.
A reusable prediction method accepts donor outcomes without requiring
target outcomes. An ATT call does not silently refit or retune.
Endpoint retention on a D-tuned fit is not independently tuned endpoint
retention; retain the tuning and prediction continuation labels separately.

The eight primary empirical cases remain: two ADH-type static cases
(uniform and empirical constant metrics), followed by three metric modes
(uniform, empirical constant, empirical time-varying) for each of $q=1$
and $q=2$. All six dynamic cases use projected increments. Within each
order-comparison pair, predictor inputs and evaluation dates are shared.
Training fit, held-out pre-treatment prediction and post-treatment ATT
are separate outputs. No Monte Carlo performance study is planned.

\section{Revised calibration and public responsibilities}\label{sec:calibration-revision}
The corresponding mathematical specification is in the revised-design
subsection of \path{tvsc-workflow-paper.tex}. The following decisions are
agreed, not claims of implemented or independently validated behavior:
\begin{itemize}
\item Combine raw block construction with \texttt{tvsc.dataprep()}.
Target and donors remain researcher-declared. Retain raw prefix access
and separate future outcomes; do not learn preprocessing in preparation.
\item Retain the existing preparation arguments and default regular step,
as listed in the approved interface in \path{tvsc-workflow-paper.tex}.
The explicit \texttt{predictors} argument lists all matching
features; do not add the outcome silently. Selecting only the outcome
column prepares the approved outcome-only specification. The intervention
index is the first treated period; training dates are explicit and
regularly spaced. Preserve supplied donor and predictor ordering and
existing training validation rules.
\item Retain class \texttt{tvsc\_prepared}, schema metadata, the ordered
raw \texttt{training} table, separate \texttt{future\_donors} and
\texttt{future\_target}, \texttt{recipe} and \texttt{diagnostics}.
Add \texttt{blocks} containing full-pre-period unscaled \texttt{X1},
\texttt{X0}, \texttt{Y1}, \texttt{Y0}, dates and dimensions. Preserve
date-indexed target feature vectors and predictor-by-donor matrices,
the target outcome vector and the date-by-donor outcome matrix. Retain
raw prefix access, unapplied-scaling metadata and constant-feature flags.
Future target outcomes remain optional. No separate public block-builder
call is required. These responsibilities are approved, not implemented;
integration validation remains pending. Object versions are specified below.
\item Version object types separately: revised preparation, raw and scaled
blocks, and scaling metadata use \texttt{schema\_version = 2L}; the new
\texttt{tvsc\_vmetrics} class starts at \texttt{1L}. Keep split objects
at version \texttt{1L} because their date-based output format is unchanged.
Versions identify object contracts, not package releases or validation.
\item Require revised preparation for metric construction and update
split-input compatibility to accept it. Do not silently upgrade old
objects. Reject incompatible preparation objects with instructions to
rerun preparation from the original panel; changing a version field
manually is not migration. Check required fields, dimensions, ordering
and relevant metadata as well as class and version.
\item The first implementation step integrates raw block construction
into preparation, updates split-input compatibility, and updates
preparation and split tests. Metric construction is a separate subsequent
step. The donor solver is outside this refactor. These versioning and
compatibility rules are now implemented in preparation and split source,
with revised fixtures passing in the supplied user-run R transcript;
independent reproduction remains pending. The separate metric-construction
source now implements the scaling and metric-object versions; its R
fixtures have a reported pass supported by the final excerpt, with the
complete log and independent reproduction still outstanding.
\item Use \texttt{tvsc.vmetrics()} for scaling and metric construction
on one actual training sample. Return transformed blocks as well as $V$,
parameters and diagnostics. Retain internal helpers where useful rather
than duplicate scaling or block logic.
\item The approved core metric interface is
\begin{verbatim}
tvsc.vmetrics(prepared, cutoff, predictor_weights = "uniform",
              scaling = "sd", conversion = "absolute",
              scales = NULL, provenance = NULL, rank_tol = NULL,
              weight_sum_tol = 1e-8)
\end{verbatim}
Require an actual pre-period date as \texttt{cutoff} and use only its
training prefix; no split object is accepted. The character weighting
inputs are exactly \texttt{"uniform"}, \texttt{"constant"} and
\texttt{"varying"}, meaning equal weights, common-slope donor OLS with
date intercepts, and separate donor regressions by date with final-date
metric retention. Do not use the longer empirical-prefixed input names.
Alternatively accept a named supplied-weight vector or predictor-by-date
matrix, aligned to the selected features and exactly the prefix dates.
Weights refer to scaled feature coordinates; no silent normalization or
slicing of extra dates is permitted.
\item Conversion accepts \texttt{"absolute"} or \texttt{"squared"}
for empirical modes and is recorded as not applicable otherwise.
Require \texttt{scales} for custom scaling and \texttt{NULL} otherwise.
Outcome-only inputs use \texttt{"uniform"} with $V_t=[1]$ and recorded
calibration bypass; reject explicit \texttt{"constant"} or
\texttt{"varying"} requests rather than silently ignoring them.
\item Return class \texttt{tvsc\_vmetrics} with schema/version metadata,
\texttt{periods}, \texttt{cutoff}, \texttt{blocks}, \texttt{weights},
\texttt{scaling}, \texttt{calibration}, \texttt{diagnostics} and
\texttt{provenance}. Blocks contain transformed predictors and unchanged
training outcomes; weights are a named $K\times T_c$ matrix. Record
scaling centers, divisors, reference donors and fallbacks, and calibration
mode, slopes, intercepts, dates, conversion and boundary convention or
the reason for bypass. Supplied provenance is not leakage certification.
No fitted donor weights, selected penalty or validation scores belong in
this object. Reuse it across orders and penalties for the same prefix and
configuration. The rank and supplied-weight sum controls below are
approved, as are the object versions above. The source draft now implements
them; R validation is pending.
\item For supplied weights, use \texttt{weight\_sum\_tol = 1e-8} as
the default absolute column-sum tolerance. Require finite nonnegative
entries, positive finite computed column sums, and
$|\sum_k v_{k,t}-1|\leq\texttt{weight\_sum\_tol}$ for every column.
Reject the supplied input if any column fails; reject even small negative
entries rather than clipping them. Preserve accepted values without
renormalization. Overrides must be finite numeric scalars in $[0,1)$;
zero requests exact equality of the computed sum to one. Record the
tolerance, column sums and maximum absolute deviation from one.
Acceptance is numerical, not exact mathematical normalization, and is
distinct from the intentional conversion of OLS slopes into normalized
empirical weights. The source implements this check; R validation remains
pending.
\item Check numerical OLS rank by SVD after accounting for date
intercepts. For \texttt{"constant"}, center donor features within each
calibration date and stack the donor rows; for \texttt{"varying"}, check
each centered donor design separately. For a design $B$ with $n$ rows and
$K$ columns, \texttt{rank\_tol = NULL} selects the relative tolerance
$\rho=\max(n,K)\epsilon_{\rm machine}$. An override is a finite numeric
scalar strictly between zero and one. Count singular values strictly
greater than $\rho\sigma_{\max}(B)$ and require rank $K$. An all-zero
design fails. Centering handles the agreed intercepts, not another
standardization; do not change matching blocks or silently rescale columns.
\item Record requested and actual rank tolerance, threshold, singular
values, numerical and required ranks, calibration dates and a
condition-number diagnostic, without an additional condition-number
rejection threshold. Declare overrides before validation and never adjust
them automatically until a fold passes. Rank rejection means calibration
failure, not dropped predictors, reduced-rank coefficients, ridge or
uniform weights. Passing this roundoff-based check may admit nearly
dependent designs and is neither a guarantee of stable predictor weights
nor certification of exact mathematical rank. The new source implements
this rule; R validation remains pending.
\item Let \texttt{tvsc.fit()} construct or accept chronological splits.
Keep \texttt{tvsc.split()} separately available; do not put validation
design into \texttt{tvsc.vmetrics()}. Fixed-$\lambda$ fitting needs no
split. Reconstruct preprocessing and empirical metrics within each prefix
and again for the final fit, before donor fitting.
\item Across uniform, empirical constant, empirical time-varying,
supplied-weight and outcome-only specifications, learn transformations
from donors only within the actual training prefix and apply them
unchanged to target and donors. No target-inclusive reference option is
planned for the first version. This is a common preprocessing rule, not
a claim of superior counterfactual accuracy. Preserve the
four scaling methods and \texttt{sd} default, with donor sample variance
denominator $J-1$. The \texttt{none} method leaves features unchanged;
\texttt{custom} applies supplied positive scales identically to donors
and target. Neither estimates reference parameters. Reconstruct learned
preprocessing within every validation prefix and the final fit.
Reporting outcomes remain unchanged. Existing
target-inclusive formulas and zero-contrast claims below are historical.
\item For genuine zero donor variation, use divisor 1 with a visible
diagnostic identifying affected features and dates. For \texttt{sd},
fallback requires donor equality within every training date for the
feature, not equality of donor levels across dates; retain no centering.
One zero-variation date does not trigger fallback if the common RMS scale
is positive. For \texttt{zscore}, handle each feature/date separately,
retaining the donor center and applying divisor 1 to both target and
donors. Preserve any target--donor gap in original feature units rather
than forcing it to zero. This is not unit-variance standardization.
Keep \texttt{none} unchanged and reject nonpositive or nonfinite
\texttt{custom} scales rather than replacing them.
\item Reserve the fallback for actual donor equality. Unequal donor
values producing a zero or nonfinite scale cause numerical failure;
do not silently floor small positive scales. Reject nonfinite transformed
values. Scaling success does not resolve unidentified OLS slopes:
donor-constant features at required calibration dates cause time-varying
calibration failure, while constant calibration may use variation at
other calibration dates if the pooled design has full rank after date
intercepts. Uniform, supplied-weight and outcome-only modes have no OLS
rank requirement from this scaling rule. The new internal helper implements
this convention; R validation remains pending.
\item Calibrate future-outcome relevance with donor-only OLS and
one-period-ahead responses wholly inside the training prefix. No ridge,
target calibration observations or post-cutoff responses are permitted.
Selected current or earlier outcomes may be features, but not the response
being predicted. Use the same scaled features in calibration and matching.
\item Retain uniform weighting as the overall default. Empirical
conversion uses normalized absolute slopes by default, normalized squared
slopes optionally, and excludes intercepts. No LMG, permutation or other
importance allocations are in implementation scope. The absolute-rule
bound argument does not establish superiority in counterfactual accuracy.
\item Use date-specific intercepts and common slopes for constant
calibration, and separate donor cross-sectional OLS regressions with
date-specific intercepts and slopes for time-varying calibration.
Intercepts enter calibration only, not predictor weights or counterfactual
construction. Require unique pooled slopes after accounting for date
intercepts in the constant mode. No rolling
windows, slope trends or temporal coefficient penalty are chosen.
For a prefix indexed $1,\ldots,T_c$, estimate the latter using pairs at
$t,t+1$ for $t<T_c$ and retain $V_{T_c}=V_{T_c-1}$. This final-date metric
retention is a boundary convention, not an equality constraint on donor
weights or a post-treatment donor-weight continuation rule.
\item Fail explicitly for all-zero or non-unique slopes. Retain individual
uniquely estimated zero weights and normalize the remaining slopes. Do
not filter by significance or an arbitrary coefficient-size threshold,
drop predictors automatically, add ridge or substitute uniform weights.
Report a failed required fold rather than omit or shift it. Final-date
retention cannot hide a failed preceding calibration regression.
\item The approved outcome-only specification uses one contemporaneous
outcome feature per training date, without covariates or an additional
lag stack. Thus $K=1$ and $V_t=[1]$. All training dates enter matching,
each with its own donor vector connected by the temporal penalty.
Bypass OLS calibration by design, not as a fallback after failure.
\texttt{tvsc.vmetrics()} still applies the chosen scaling and returns
the one-feature metric. Retain the four scaling methods, \texttt{sd}
default, both penalty orders, chronological tuning and continuation rules.
At fixed $\lambda$, scaling changes the matching--regularization balance;
reporting outcomes remain unchanged. Within validation, fit only on the
training prefix and reserve subsequent target outcomes for scoring.
Preparation and metric construction now reflect this specification in
source; fitting integration and R validation of metric construction remain
pending.
\end{itemize}

Fitting arguments and schemas remain open. The approved object-versioning
and compatibility rules require implementation and validation, as do the
approved numerical rank and
supplied-weight sum rules require implementation and validation.
The approved preparation and metric
contracts and zero-scale convention must be
implemented and tested separately from OLS rank failures.
The empirical comparison layout must
also be reconciled with the additional outcome-only specification.

\section{Primary-documentation findings and consequences}
\paragraph{Static data preparation is not date-specific preparation.}
The Synth manual separates predictor summaries, fitting-outcome matrices
and plotting outcomes. Its \texttt{dataprep()} supports aggregation over
declared predictor periods and separate special predictors; the manual
also documents missing-value removal in predictor aggregation
\cite{synth}. Our design consequence is to borrow the separation of
roles, not silently adopt aggregation or missing-data rules for changing
weights. A repeated pre-period average is not a contemporaneous block.

\paragraph{Scaling has an orientation and a training sample.}
Base R's \texttt{scale()} works on columns. With centering enabled, it
uses sample standard deviations; without centering, its default scaling
uses root mean squares. It retains the fitted centers and scales as
attributes \cite{rscale}. Since our $X_{0,t}$ has predictors in rows and
donors in columns, blindly calling \texttt{scale(X0)} would scale donors
rather than predictor coordinates. The proposed solution is to estimate
predictor-level transformations on a long training table, then construct
the matrices and preserve the transformation metadata.

\paragraph{Preprocessing should be estimated and then applied.}
The recipes documentation distinguishes estimating transformations from
training data and applying them to other data \cite{recipes}. We propose
the same separation internally, without requiring recipes as a dependency:
a raw prepared object plus a cutoff-specific block builder. Reusing one
full-pre-period standardized matrix inside shorter validation prefixes is
not the intended workflow.

\paragraph{Chronological windows must respect the panel's date groups.}
The rsample documentation distinguishes row-based sliding windows from
index-based windows and permits expanding windows \cite{rsample}.
Our design consequence is to construct folds on unique ordered dates,
then select every eligible unit at those dates. Random row splits and
row-count windows over a long panel are not substitutes for this operation.

\paragraph{The primary quadratic form can be singular.}
The quadprog manual describes strictly convex quadratic programs, whereas
OSQP documents positive-semidefinite quadratic objectives
\cite{quadprog,osqp}. Thus OSQP is a candidate backend to investigate,
not a selected or benchmarked dependency. Do not add a small ridge merely
to satisfy a solver's requirements: that changes the estimator. Numerical
primary and secondary acceptance require explicit tests.

\paragraph{The ADH comparison can accept supplied predictor weights.}
The Synth manual documents \texttt{custom.v} as bypassing its predictor-
weight optimization \cite{synth}. This provides a possible adapter for
our uniform and externally calibrated static cases. Its predictor
construction and effective scaling must be checked before use; avoiding
double standardization is a release task. Neither modified case is
automatically a reproduction of the published estimates.

\section{Raw preparation implementation and future contract}\label{sec:prep}
The revised interface in \path{R/tvsc.dataprep.R} now adds version-2
raw blocks. Its revised fixtures pass in the supplied September 15, 2026
user-run R transcript, including isolated source loading. Independent
reproduction remains pending; the earlier user-run pass concerns the
preceding raw-panel-only draft. The argument list is
unchanged:
\begin{verbatim}
tvsc.dataprep(
  data, unit, time, outcome, treated, donors, intervention_time,
  pre_period, predictors, period_step = 1
)
\end{verbatim}
Column roles are supplied as names, not positional column numbers.
Predictor names must currently be explicit. Including the outcome's column
name includes it once; omitting it excludes it from matching while retaining
it for scoring. This avoids settling the proposed automatic outcome-inclusion
default at this stage. Scaling has no public argument yet: the stored
\texttt{recipe\$scaling = NULL} marks an unresolved operation, not an
approved unscaled fitting procedure.
The function validates and stores raw inputs and recipe metadata;
it estimates no donor weights, chooses no regularization parameter and
constructs no validation splits. The proposed \texttt{intervention\_time}
denotes the first treated period; \texttt{pre\_period} must precede it.
That eligibility boundary is distinct from optional splitting within the
pre-period. A future scaling-recipe interface still needs to be finalized.
No reporting horizon is required here. Actual preprocessing parameters
will be estimated later, when the fitting routine defines its full training
sample or a validation prefix.
Any full-training matrix preview must be labeled as such and must not be
used as a shortcut for prefix-specific preparation.

The revised object has class \texttt{tvsc\_prepared}, schema version 2,
and fields \texttt{schema}, \texttt{training}, \texttt{future\_donors},
\texttt{future\_target}, \texttt{recipe}, \texttt{diagnostics} and
\texttt{blocks}. The blocks hold the full raw training arrays with their
own version-2 schema and preserve dimensions when there is one feature.
Training rows are ordered by declared date, then target followed by donors
in user-specified order. Original measurement values are retained. Constant
predictors are recorded but not removed, since no scaling is attempted.
Only required training cells receive completeness and finite-value checks;
future outcomes are stored with an explicit unvalidated flag. Their
duplicate keys, missingness and prediction eligibility must be checked by
the later prediction/scoring layer. Schema-level unit and time validation
still applies to all supplied rows.

Example calls, once the R file is sourced from the project directory:
\begin{verbatim}
source("R/tvsc.dataprep.R")
prepared <- tvsc.dataprep(
  data = my_panel, unit = "state", time = "year", outcome = "sales",
  treated = "Target", donors = c("DonorA", "DonorB"),
  intervention_time = 2003, pre_period = 2000:2002,
  predictors = c("sales", "income")
)
\end{verbatim}
This example specifies a hypothetical dataset, not California results.
The fixture file creates its own small dataset. Run it with
\texttt{Rscript tests/test-tvsc.dataprep.R} from the project root when R is
available. Its cases include row permutations, preserved donor order,
duplicate/missing training keys, nonfinite values, explicit predictor
inclusion, absent future target outcomes, constant predictors and changes
to future numerical values that must not change prepared training values.
Until that command succeeds, this is an unvalidated code draft.

\subsection{Panel structure and block definition}
\begin{enumerate}
\item Require one target and at least two explicitly named eligible
donors. Preserve the supplied donor order and original identifiers.
Reject a target appearing in its own donor set and duplicate donor IDs.
\item The current draft uses integer-valued numeric periods and an explicit
\texttt{period\_step} (default 1), with at least three training dates.
The declared pre-period must be increasing at that step. Retain actual
labels separately. Do not compress missing calendar periods into adjacent
observations. Monthly calendar labels require a proper period index, not
an assumption that months have equal numbers of days.
\item Require unique unit--date keys and finite numeric outcomes and
predictors in the required training cells. Report offending keys rather
than silently aggregate, impute, drop donors or change the training window.
Validate future prediction cells when they are actually requested.
\item Contemporaneous blocks are the intended first construction; the
current code retains their raw columns without building matrices.
The proposed \texttt{include\_outcome} convenience option remains deferred;
the current implementation requires explicit predictor names and rejects
duplicate names. Outcome-only preparation uses that outcome column alone.
The response used later for empirical calibration must not be regressed
on an identical copy of itself.
\item Preserve time-invariant characteristics that differ across units.
Do not remove a predictor simply because it is constant within each unit.
Reject automatic factor-to-integer conversion; categorical encoding needs
an explicit later contract.
\item Defer lags and trailing summaries until their history requirements,
first usable date and prefix reconstruction are specified. Never use a
future target observation to create a held-out prediction feature.
\end{enumerate}

\subsection{Economic scaling: agreed initial scope}
This subsection records the existing target-inclusive source draft.
For all revised weighting modes, the donor-only learned-scaling decision in
Section~\ref{sec:calibration-revision} supersedes this reference population.
The new internal scaling helper in \path{R/tvsc.vmetrics.R} implements
donor-only scaling and revised zero-scale rules. A reported fixture pass
is supported by the final excerpt; full-log review and independent
reproduction remain pending. The legacy standalone source described here is unchanged.
On September 15, 2026 the user agreed to short panel-specific names and
to defer the other methods. The initial options are \texttt{sd} (default),
\texttt{zscore}, \texttt{none} and \texttt{custom}. A scaling source draft
and fixtures are now written and the displayed checks passed in the fourth
user-supplied transcript, but have not been run in R here. Existing preparation and block assembly remain
unchanged. Let $N=J+1$ include the target and donors, and let $T_c$ be the
training-prefix length. Define
\[
 m_{k,t}=\frac1N\sum_{i=1}^N X_{i,k,t},\qquad
 s_{k,t}^2=\frac1{N-1}\sum_{i=1}^N(X_{i,k,t}-m_{k,t})^2,
 \qquad s_{k,c}=\sqrt{T_c^{-1}\sum_{t=1}^{T_c}s_{k,t}^2}.
\]
The \texttt{sd} option uses $X_{i,k,t}/s_{k,c}$ without centering;
\texttt{zscore} uses $(X_{i,k,t}-m_{k,t})/s_{k,t}$ separately by date.
Thus \texttt{sd} uses an RMS within-date SD, not pooled-level SD, whereas
\texttt{zscore} changes the divisor across dates. These are not merely
centered and uncentered versions of the same matching criterion.
\texttt{none} preserves all supplied values. \texttt{custom} divides by
finite positive scales matched by predictor name and common across dates.
Pooled-level SD/z-scoring, MAD, IQR, min--max, maximum-absolute-value and
other robust options are outside the initial scope. Do not silently add them.
No unit-specific transformation or post-cutoff observations enter these
estimates. Preserve supplied reporting outcomes unchanged. Future prediction
does not use target-based z-scores. Store the method's precise definition,
reference units, dates, centers, scales and diagnostic decisions.

Under unit-sum donor weights, common centering cancels in each matching
discrepancy. Algebraically,
\[
 \frac{X_{1,k,t}-m_{k,t}}{s_{k,c}}
 -\sum_jw_{j,t}\frac{X_{0,k,j,t}-m_{k,t}}{s_{k,c}}
 =\frac{X_{1,k,t}-\sum_jw_{j,t}X_{0,k,j,t}}{s_{k,c}}.
\]
Scaling does not cancel. This identity motivates an implementation test,
not a claim that all preprocessing conventions are equivalent.
Writing the original-unit discrepancy as $e_{k,t}$ gives
\[
 \sum_kv_{k,t}(e_{k,t}/s_{k,c})^2
 =\sum_k\frac{v_{k,t}}{s_{k,c}^2}e_{k,t}^2.
\]
Thus uniform $v$ means equal coefficients in scaled units, not equal
economic importance. For \texttt{zscore}, replace $s_{k,c}$ by $s_{k,t}$;
the effective coefficients then change by date. Define the economic variables first, including their
real versus nominal, per-capita or rate interpretation. Then choose the
matching yardstick and document it separately from predictor weights.

If $X_{i,k,t}=a_{k,t}+z_{i,k,t}$, the common component $a_{k,t}$ cancels
from the matching discrepancy under unit-sum weights, but can increase
a pooled standard deviation across time. Pooled scaling can consequently
downweight discrepancies because of a common trend absent from those
discrepancies themselves. Within-date dispersion avoids that particular
effect but remains sensitive to donor composition and small dispersion.
This is the rationale for the agreed default, not proof of optimality or
statistical consistency. Conditional results apply to finite transformed
blocks, not an unstated asymptotic calibration or tuning guarantee. The
preparation layer preserves raw values and does not itself apply the default.

The user approved exact-equality scale-1 fallbacks: under \texttt{sd},
every training date must have identical values across all reference units,
although the predictor may vary over time; under \texttt{zscore}, an
individual equal-valued date uses its common value as center and divisor
1, producing zeros. Flag the affected predictor-date cells. A zero or
nonfinite computed SD despite unequal values is an error. Positive scales
are not clipped or floored; nonfinite transformed values are rejected.
Record the smallest applied divisor without claiming good conditioning.
Do not drop predictors, change $K$ or renormalize $V$.

The implemented draft signature is
\begin{verbatim}
tvsc.scale(blocks, scaling = "sd", scales = NULL)
\end{verbatim}
It accepts raw version-1 blocks, validates metadata and array alignment,
and returns a new object of class
\texttt{c("tvsc\_scaled\_blocks", "tvsc\_blocks")}. Outcomes, input
schema and original raw diagnostics are unchanged. \texttt{none} performs
no arithmetic on predictor arrays. Every method records a non-null
\texttt{recipe\$scaling} and rejects repeated processing, including after
\texttt{none}. Custom scales must be finite, positive and named exactly
once per predictor; align by name and reject non-null scales for other
methods. No function dependencies beyond base R and its stats namespace
are needed by the scaler itself.

The recipe records its schema version, method, formula identifier,
reference units, training periods and cutoff, applied centers and divisors
as predictor-by-date matrices, estimated SDs where applicable, exact
zero-variation and fallback masks, a zero-policy label, and diagnostics.
None/custom have null estimated SDs and a not-estimated reference mode.
The RMS calculation normalizes by the largest within-date SD before
squaring to avoid unnecessary overflow. Rebuild learned transformations
in every training prefix. R fixtures cover formulas, identities, outcomes,
metadata, cutoff isolation, invalid input and numerical failures. The
displayed checks passed in the fourth user-supplied transcript; independent
reproduction and the standalone example remain pending. No R executable is available
in this workspace and no dependencies were installed.

\subsection{Objects and information boundaries}
Propose a classed prepared object with logically separated raw inputs and
recipes; realized transformations belong to training-builder outputs:
\begin{itemize}
\item Raw eligible pre-treatment data and a schema: IDs, dates, intervention
boundary, predictor names and order, outcome inclusion, units and block recipe.
\item A scaling recipe with no fitted parameters at preparation time.
Training-specific centers, scales and diagnostics are generated by an
internal builder and stored with the resulting training view and fitted
model, not reused across incompatible prefixes.
\item Future donor outcomes, if supplied, isolated from training and
validated by prediction for its requested dates and donor identities.
\item Observed future target outcomes, if supplied, reserved for the
ATT/scoring operation. Their absence must not prevent fitting or donor-
conditional prediction; future target predictors are unnecessary.
\end{itemize}
The agreed separation now has a public raw assembly function,
\texttt{tvsc.blocks(prepared, cutoff)}, implemented without transformations.
This replaces the raw-assembly portion of the earlier proposed internal
\texttt{build\_training} helper; a later training layer still needs explicit
scaling and predetermined predictor metrics. The raw builder returns
$X_{1,t}$ as a length-$T_c$ list of named length-$K$ vectors, and
$X_{0,t}$ as a matching list of $K\times J$ matrices; outcome arrays in
supplied reporting units have dimensions $T_c$ and $T_c\times J$. The returned
\texttt{tvsc\_blocks} list contains \texttt{schema\_version}, a schema whose
pre-period is restricted to the prefix, \texttt{periods}, \texttt{cutoff},
\texttt{dimensions}, \texttt{X1}, \texttt{X0}, \texttt{Y1}, \texttt{Y0},
\texttt{recipe} and \texttt{diagnostics}. All values remain as supplied;
\texttt{scaling = NULL} records no package scaling, not upstream provenance
or an approved default for the future estimator.
Constant predictors are flagged using the prefix only, not the full panel.
The cutoff is required and must be an actual declared pre-period date
leaving at least three dates. The builder checks schema/time metadata,
then reuses \texttt{tvsc.dataprep()} on prefix rows to validate keys and
required numeric values and restore target/donor ordering. Later values and
future tables are not passed to that validator or included in the output.
There are no lags, summaries, imputations or fitted transformations.
The public builder therefore requires \texttt{tvsc.dataprep()} to be loaded.
Never
rely on incidental input row order or matrix dimension dropping at $K=1$.
The fitting routine receives only the training view. With no validation,
the builder uses the full designated pre-period; with validation it is
called separately for each training prefix. These boundaries are contracts
to be tested, not a claim that an R object provides a security
sandbox against arbitrary user-written callbacks.

\section{Separate chronological split contract}\label{sec:splits}
The agreed responsibility of \texttt{tvsc.split()} is to create date
assignments and fold metadata, not to prepare transformed copies of the
data or perform fitting. The agreed call is
\begin{verbatim}
splits <- tvsc.split(prepared, initial, horizon, step = 1)
\end{verbatim}
Expanding training prefixes, explicit initial and horizon arguments and
\texttt{step = 1} are agreed. Window arguments count
regular periods, not long-panel rows. Each fold records its ordered
training dates, fitting cutoff, subsequent validation dates and horizons,
with all selected units included at every assigned date. Require at least
three training dates, a nonempty later validation window, positive integer
window arguments, no overlap between training and validation within a
fold, and dates confined to the eligible pre-period. Preserve the period
index and donor/schema identifiers for downstream compatibility checks.
The current implementation retains full validation windows only: cutoffs
advance by the requested stride, with no shortened final window or extra
off-stride cutoff. If no full fold fits, it raises an error. This boundary
policy is explicit in the documentation and returned specification.
With eight dates, initial size four and horizon two, the three folds are
1--4/5--6, 1--5/6--7 and 1--6/7--8.

The returned \texttt{tvsc\_splits} object contains a schema version, a copy
of the input schema, the specification, ordered folds and diagnostics.
Each fold has its identifier, date positions within the declared pre-period,
actual training and validation dates, cutoff and period-count horizons.
Diagnostics include fold count, assessment counts by date and unassessed
dates, including initial training-only dates and stride/tail omissions.
Repeated assessment dates are exposed without selecting scoring weights.
No independent test set is reserved. The constructor uses schema metadata
only and does not revalidate raw values. A schema copy is not a data checksum;
downstream code must validate compatibility rather than assume identical
data values merely because schemas match.

\texttt{tvsc.fit()} checks the supplied split specification when
\texttt{lambda = "cv"}. With window sizes deliberately explicit, a missing
specification requires an informative request/error rather than an
invented split or substitution of training loss. A fixed numeric penalty
does not require splitting or out-of-sample scoring. Held-out assessment
can use the same date-assignment machinery but must remain separate from
tuning; no function name makes reused tuning outcomes an independent test.
The constructor itself learns no preprocessing parameters, predictor
metrics or donor weights and reads no outcome values to choose windows.

\section{Predetermined metrics and chronological validation}
The initial \texttt{tvsc.metrics()} draft accepts exactly
\texttt{"uniform"}, a named common real numeric vector or a real numeric
predictor-by-date matrix with explicit row and column names. It requires
schema-version-1 \texttt{tvsc\_scaled\_blocks}; choosing
\texttt{tvsc.scale(..., "none")} makes no-scaling explicit. It aligns
names without partial matching and rejects missing or extra predictors
or dates rather than silently slicing. It never calibrates weights.

For a prefix of length $T_c$, uniform inputs give $v_{k,t}=1/K$, a common
vector gives $v_{k,t}=a_k$, and a matrix gives $v_{k,t}=M_{k,t}$. Entries
must be finite and nonnegative, with positive column sums satisfying
$|\sum_kv_{k,t}-1|\leq\varepsilon$. The recorded absolute tolerance
defaults to $10^{-8}$ and must satisfy $0\leq\varepsilon<1$.
No clipping or renormalization is applied; an accepted near-unit sum is
not asserted to be exactly one. Zero weights are allowed.

The output is a \texttt{tvsc\_metrics} object containing the canonical
$K\times T_c$ matrix, shape-based input mode, schema, cutoff, scaling
record and diagnostics. It does not copy or read numerical predictor or
outcome arrays. Metadata alignment is checked without certifying those
arrays or recalculating scaling; the eventual fitter must validate actual
numerical inputs and their compatibility. A supplied matrix retains mode
\texttt{supplied\_time\_varying} even if its columns happen to be equal;
exact column constancy is a separate diagnostic. No input shape proves
empirical origin. Optional provenance is a character description stored
with \texttt{verified = FALSE}, not a leakage certificate. The constructor
does not renormalize effective raw-unit coefficients after scaling.

For cross-validation, distinguish genuinely external supplied metrics
from empirical metrics estimated from this panel. Merely subsetting a
full-pre-period empirical matrix does not establish absence of leakage:
its entries may have used later target outcomes. Empirical modes need a
calibration callback or fold-specific metric objects with declared training
provenance. Each callback receives only the allowed prefix, runs before
donor fitting and is not reoptimized for each $q$ or $\lambda$.
The final full-training calibration is also frozen before the final fit.
The donor-only OLS direction and coefficient conversions are now agreed
in Section~\ref{sec:calibration-revision}. Remaining model and numerical
contracts listed there must be resolved before replacing this constructor
with the combined \texttt{tvsc.vmetrics()} operation.

The complete chronological selection operation is:
\begin{enumerate}
\item Declare assessment windows, tuning cutoffs, initial prefix length,
validation horizons, stride and any boundary-expansion rule before scoring.
Use folds defined on dates and include all selected units in each fold.
\item For each tuning prefix, rebuild blocks and scaling, then calibrate
or validate its supplied metric. Cache these fixed inputs for both orders
and all candidate penalties within that metric mode.
\item Fit each candidate with the specified secondary selection. Extend
the final fitted increment from that prefix, using contemporaneous donor
outcomes but not target validation outcomes.
\item Reveal target validation outcomes for scoring only. Use the paper's
mean of fold-specific mean squared errors. Record horizons and per-date
errors; overlapping windows are not independent replications.
\item Reject candidates with any required fit failure; never compare
averages over candidate-specific successful subsets. If no candidate is
eligible, return selection failure. Preserve the original grid and statuses.
\item Apply a predeclared tie rule. Proposed rule: choose the largest
$\lambda$ among exactly equal minimum scores; do not call this a statistical
one-standard-error rule. Any near-tie tolerance requires a separate decision.
\item Assess the selected procedure on untouched later pre-treatment dates
when feasible, then refit on all designated pre-treatment dates. If an
assessment cutoff is used for nested evaluation, all tuning for that cutoff
must occur within its own earlier prefix.
\end{enumerate}
Initial-window and horizon sizes are deliberately explicit, not defaulted.
A short-horizon validation result must not be described as verification
of a much longer post-treatment extrapolation. In particular, a short
pre-period may not support long-horizon tuning and an independent test.

\section{Computation, prediction and reporting}
\paragraph{Computation proposal.}
Use sparse block matrices and construct the difference operators once per
prefix and order. Investigate ordered-grid warm starts. OSQP documents
warm starts and factorization reuse for suitable parameterized problems
\cite{osqp}; changing $\lambda$ changes the quadratic matrix, so numerical
factorization reuse must not be promised for every candidate. Cache
preprocessing and sparsity patterns separately from numerical factorizations.
Pin and test an actual R interface version before adopting its API.

With $F$ tuning folds, the six dynamic cases imply $6\cdot100\cdot F$
candidate fits before final refits, outer assessments, static cases or
secondary solves. This is a workload count, not a runtime estimate.
At $\lambda=0$, the two orders have the same fitted objective and selection,
so an exact shared result may be cached while retaining both method labels.
Secondary selection must be addressed during tuning too, because different
primary optima can have different endpoints and predictions. Warm-start
choice must not silently replace minimum-norm selection.

\paragraph{Prediction and ATT contract.}
For each requested period use its horizon relative to the training cutoff,
not its position in a filtered output table:
\[
 \widehat w^E_{T+h}=\widehat w_T,\qquad
 \widehat w^D_{T+h}=\Pi_{\Delta_J}
   \{\widehat w_T+h(\widehat w_T-\widehat w_{T-1})\}.
\]
Projection is Euclidean simplex projection, not clipping followed by
renormalization, and not recursive projection. Return projected weights,
projection distances and original-unit donor-weighted predictions.
Only then join target outcomes by date to form gaps and their average on
the declared complete window. Missing required donor outcomes cannot be
handled by silently redistributing their weights. Missing target outcomes
may leave predictions available while making the declared ATT unavailable.
Save training and prediction configurations so E/D changes are transparent.

\paragraph{Static benchmark adapter.}
Keep the ADH summaries and fitting periods explicit and distinct from
our date-specific blocks. Inspect the chosen Synth version's scaling and
\texttt{custom.v} behavior before passing inputs; capture effective metrics
and settings. Original ADH reproduction may use its original metric-selection
procedure as a separately labeled benchmark, without changing our own
predetermined-metric estimator. Do not present agreement with ADH as truth.

\section{Implementation sequence and acceptance gates}
The sequence below is retained as an earlier roadmap, not the current
restart order. Settle and refactor the revised preparation/calibration
contracts first. The \texttt{tvsc.problem(blocks, metrics, lambda, q = 1)}
draft is set aside pending that work. It revalidates supplied metric
metadata and canonical weights, checks numerical predictor arrays, and
returns the dense primary matrices $P,c$, objective constant
$(1-\lambda)b^\top b$, normalized matching/difference operators, date-wise
simplex constraints and a date/donor coordinate map. Reporting outcomes
are neither read nor numerically validated here. It accepts singular
objectives without a ridge and rejects nonfinite arithmetic. There is no
solver, numerical solution acceptance or secondary selection. The
unweighted stored difference operator must be multiplied by
$\sqrt\lambda$ for any later secondary penalty-image condition.
Its fixtures compare the matrix objective with direct losses and check
both orders' null spaces and zero-penalty equivalence. Dense reference
storage is not a production memory or runtime claim. Do not choose a
solver backend merely to avoid implementing the stated estimator.

\begin{enumerate}
\item Confirm the preparation decisions below. Implement a small pure
preparation module and prefix builder before adding a solver or calibration.
Gate: named matrices, original outcomes and prefix metadata match hand-
calculated fixtures in R.
\item Implement the separate split constructor. Gate: date-level membership,
ordering, intervention eligibility, minimum lengths and compatibility checks
pass. The preparation object remains usable without constructing splits.
\item Implement both fitting orders, numerical statuses and minimum-norm
selection. Gate: simplex constraints, objective scaling and known small
solutions pass, including singular designs and $\lambda=0$.
\item Implement projection and donor-conditional prediction. Gate: known
simplex projections, nonconsecutive requested horizons and both continuation
rules pass without requiring target outcomes.
\item Implement chronological tuning. Gate: fold-specific transformations,
scores, failure handling and continuation labels reproduce saved fixtures.
\item Implement ATT and summaries. Gate: date joins, original units,
window completeness and direct averages agree exactly up to stated tolerance.
\item Specify and validate empirical metric calibration and the static
adapter. Gate: all eight declared cases are implementable with documented
inputs, without future-target leakage or hidden changes to the estimator.
\item Measure runtime in R and record numerical tolerances, package
versions and the execution environment. Complete paper two's reproducible
examples before returning to paper one for the California results.
\end{enumerate}

Preparation tests must cover permuted rows, preserved donor order,
duplicated keys, a donor equal to the target, $K=1$, absent future target
data, constant and time-invariant predictors, nonfinite required training
cells, irregular dates and common scaling. Include a no-split preparation
and direct fixed-penalty workflow, and separate tests for invalid or missing
CV splits. Check that split construction does not alter raw data or learn
transformations. Mutation tests should change
post-cutoff values without changing a prefix's prepared training inputs.
Within a tuning fold, changing validation target values may change scores
and ultimately the selected penalty, but must not change that fold's
candidate fits or their predictions. These tests are software validation,
not a Monte Carlo comparison of causal estimators.

\section{Decisions to confirm and exact restart point}\label{sec:decisions}
\begin{enumerate}
\item Confirm contemporaneous numeric predictors with the outcome included
once by default, plus an explicit option to exclude it.
\item Preserve the agreed four-option scaling scope and panel-specific
\texttt{sd} default and approved zero-scale/metadata contract. Execute
\path{tests/test-tvsc.scale.R} and the standalone example in a clean R
session before claiming validation. Pooled-level and robust alternatives
remain deferred.
\item Confirm the minimal numeric period-index interface and preservation
of user-specified donor ordering.
\item The separate, optional \texttt{tvsc.split()} responsibility is agreed.
Expanding windows and explicit sizes are agreed; the draft uses full windows
only. Its fixtures passed in the user-supplied transcript; reproduction and
publication metadata remain pending. Settle subsequent
tuning tie-breaking, along with
the empirical metric calibration models, solver/backend versions and
numerical acceptance settings. These do not block designing the raw panel
schema, but they do block claiming a complete implemented procedure.
\end{enumerate}
The immediate validation checkpoint is a complete clean-session log of
\path{tests/test-tvsc.vmetrics.R} and regression reruns of preparation and
splitting, recording session and source identifiers. The new metric-construction source reconstructs raw
prefixes, applies donor-only scaling once, and implements the approved
OLS, conversion, rank and supplied-weight rules. Its fixtures include
\texttt{lm.fit()} comparisons and isolation/failure checks. The reported
pass is supported only by the final isolated-loading excerpt, not a full
suite log or independent reproduction. Do not advance the solver or describe the earlier tests
as covering this revision. Remaining fitting contracts still require
agreement before implementation.
Independently, obtain the preceding GitHub handoff's integration
summary and any patches. The four earlier suites passed in the displayed
renamed user run; metrics has a subsequent user-reported pass with only
the ending supplied. Record clean-session file loading, full fixture runs,
standalone examples, source revisions and \texttt{sessionInfo()} during
integration. Do not describe those reports as independent reproduction or
as evidence for the new builder. Calibration's remaining interface and
numerical choices, solver/backend choice and numerical acceptance settings
remain open; its agreed scientific scope is recorded above.

On resumption, read this note and the corresponding section of
\path{tvsc-workflow-paper.tex}; inspect the current files and working-tree
changes rather than assume the code draft has been validated. Recheck R availability.
The split-test transcript has been reviewed and shows no failed checks.
The next publication task is to record the tested revisions and R session
information and check any remaining standalone examples. The equivalent
eight-date split example has already passed within the supplied fixtures.
The preparation publication record and standalone example remain pending
separately. Publish each completed
function to the agreed GitHub destination; publication revisions have not
been recorded here. The raw builder's displayed checks have now passed in
the user-supplied transcript. Reproduce \path{tests/test-tvsc.blocks.R} in a
clean session, run the standalone example and record the tested revision
and session information for publication. The displayed scaling checks have
passed. Hand off the scaler for GitHub integration, clean-session fixture
reproduction, standalone-example execution, prior-function regression
reruns, and revision/session records before publication. Do not add
transformations to the preparation or raw block builder implicitly.
Do not assume that the entire proposed
preparation contract has been implemented. Preserve the agreed separation,
four-option scope and panel-specific default on resumption; do not reinstate
pooled scaling as a default. Do not
resume the consistency discussion, alter paper one's estimator, execute
the California comparison or infer approval of the proposals above.

\begin{thebibliography}{9}
\bibitem{synth}
Synth authors. \emph{Synth package reference manual}, version 1.1-10,
accessed September 15, 2026. Sections \texttt{dataprep} and \texttt{synth}.
\url{https://stat.ethz.ch/CRAN/web/packages/Synth/Synth.pdf}.
\bibitem{rscale}
R Core Team. \emph{Scaling and Centering of Matrix-like Objects}.
Base R documentation, accessed September 15, 2026.
\url{https://www.stat.ethz.ch/R-manual/R-devel/library/base/html/scale.html}.
\bibitem{recipes}
recipes authors. \emph{Estimate a preprocessing recipe}.
Official documentation, accessed September 15, 2026.
\url{https://recipes.tidymodels.org/reference/prep.html}.
\bibitem{rsample}
rsample authors. \emph{Time-based Resampling}.
Official documentation, accessed September 15, 2026.
\url{https://rsample.tidymodels.org/reference/slide-resampling.html}.
\bibitem{quadprog}
quadprog authors. \emph{Functions to Solve Quadratic Programming Problems}.
Package reference manual, accessed September 15, 2026.
\url{https://stat.ethz.ch/CRAN/web/packages/quadprog/quadprog.pdf}.
\bibitem{osqp}
OSQP developers. \emph{OSQP solver documentation} and \emph{The solver}.
Official documentation, accessed September 15, 2026.
\url{https://osqp.org/docs/}; \url{https://osqp.org/docs/solver/}.
\end{thebibliography}
\end{document}